\documentclass[10pt,letterpaper]{article}
\usepackage[utf8]{inputenc}
\usepackage[T1]{fontenc}
\usepackage[margin=0.35in]{geometry}
\usepackage{xcolor}
\usepackage[hidelinks]{hyperref}

\definecolor{primary}{RGB}{18, 38, 68}
\definecolor{textdark}{RGB}{30, 30, 30}
\definecolor{linkcolor}{RGB}{0, 75, 145}

\hypersetup{colorlinks=true, urlcolor=linkcolor, linkcolor=linkcolor}
\pagestyle{empty}
\raggedbottom
\raggedright

\newcommand{\resumesection}[1]{%
  \vspace{2pt}%
  {\color{primary}\normalsize\bfseries\MakeUppercase{#1}}%
  \vspace{1pt}\hrule height 0.5pt\vspace{2pt}%
}

\def\labelitemi{\color{primary}$\bullet$}
\renewenvironment{itemize}{%
  \begin{list}{\color{primary}$\bullet$}{%
    \setlength{\leftmargin}{1.1em}%
    \setlength{\itemsep}{0.5pt}%
    \setlength{\topsep}{1pt}%
    \setlength{\parsep}{0pt}%
    \setlength{\partopsep}{0pt}%
  }%
}{%
  \end{list}%
}

\newenvironment{skilllist}{%
  \begin{list}{}{%
    \setlength{\leftmargin}{0pt}%
    \setlength{\itemsep}{0.5pt}%
    \setlength{\topsep}{1pt}%
    \setlength{\parsep}{0pt}%
    \setlength{\partopsep}{0pt}%
  }%
}{%
  \end{list}%
}

\begin{document}
\color{textdark}

\begin{center}
    {\LARGE \textbf{\color{primary} Alex Rivera}} \\ \vspace{1.5pt}
    {\large \textbf{Ingeniero de Software de Sistemas \textbar\ C++, Optimización de Rendimiento \& Concurrencia}} \\ \vspace{2pt}
    {\small
        Remote / Global \ \textbullet \
        +1 (555) 019-2834 \ \textbullet \
        \href{mailto:alex.rivera@example.com}{alex.rivera@example.com} \ \textbullet \
        \href{https://alexrivera.dev}{alexrivera.dev} \ \textbullet \
        \href{https://github.com/alexrivera-dev}{github.com/alexrivera-dev} \ \textbullet \
        \href{https://linkedin.com/in/alexrivera-dev}{linkedin.com/in/alexrivera-dev}
    }
\end{center}

\vspace{-6pt}

\resumesection{Perfil Profesional}
{\small Ingeniero de Software y estudiante de Ciencias de la Computación en la Technical University Institute con sólida formación en sistemas computacionales, programación a bajo nivel y optimización de rendimiento. Especializado en C++17, arquitectura de procesadores ARM y desarrollo concurrente en Go y Java. Experiencia comprobada en optimización de ciclos de CPU (-83\% de ciclos en hardware embebido), análisis de memoria (DMA, IWRAM/EWRAM), desarrollo de servidores de red sin frameworks y docencia universitaria en concurrencia y algoritmos.}

\resumesection{Proyectos de Sistemas y Rendimiento}

\textbf{VioletHat — Motor Gráfico de Raycasting de Alto Rendimiento} \textbar\ \textit{C++17, ARM, SDL2} \hfill \href{https://github.com/alexrivera-dev/gba_raycaster}{github.com/alexrivera-dev/gba\_raycaster}
\begin{itemize}
    \item Desarrollé desde cero un motor de renderizado DDA con texturas en C++17 para hardware bare-metal Game Boy Advance (ARM7TDMI) y sistemas de escritorio (Linux/Windows vía SDL2).
    \item Reduje el consumo de CPU de \textbf{6.46M a 1.12M ciclos de reloj por frame} (\textbf{reducción del 83\%}), elevando los cuadros por segundo de \textbf{2.6 a 14.9 FPS} sostenidos mediante profiling exhaustivo de cuellos de botella de hardware.
    \item Diseñé un pipeline de renderizado con aritmética de punto fijo, transferencias directas de memoria (DMA) en VBlank, colocación estratégica de código crítico en memoria ultrarrápida IWRAM y generación procedural de mapas.
    \item Implementé suite de pruebas automatizadas y compilación multiplataforma en GitHub Actions con artefactos ROM verificados.
\end{itemize}

\vspace{1.5pt}

\textbf{Servidor HTTP/1.1 Concurrente desde Cero} \textbar\ \textit{Go, TCP Sockets, Sincronización} \hfill \href{https://github.com/alexrivera-dev/go-http}{github.com/alexrivera-dev/go-http}
\begin{itemize}
    \item Implementé un servidor de red HTTP/1.1 completo sobre sockets TCP nativos sin dependencias externas, programando el parser formal de especificación RFC 7230, multiplexación de rutas y control de estado de conexiones.
    \item Optimicé el rendimiento concurrente utilizando un modelo de goroutines con pools de buffers reciclables para minimizar la presión sobre el Garbage Collector y maximizar el throughput de peticiones.
\end{itemize}

\vspace{1.5pt}

\textbf{Intérprete de Lenguaje LISP} \textbar\ \textit{Java, Compiladores, Estructuras de Datos} \hfill \href{https://github.com/alexrivera-dev}{github.com/alexrivera-dev}
\begin{itemize}
    \item Construí un intérprete de lenguaje funcional desde cero: análisis léxico (tokenizador), parser descendente recursivo a Árbol de Sintaxis Abstracta (AST) y motor de evaluación recursiva con entornos dinámicos de memoria y closures.
\end{itemize}

\vspace{1.5pt}

\textbf{Proyectos de Automatización con Microcontroladores STM32} \textbar\ \textit{C, ARM Cortex-M4, Hardware} \hfill Tech University
\begin{itemize}
    \item Programé firmware en C para placas STM32F446RE (ARM Cortex-M4), interactuando directamente con registros de memoria para configuración de timers de alta precisión, puertos GPIO, interrupciones externas (EXTI) y buses de comunicación.
\end{itemize}

\resumesection{Experiencia Profesional}

\textbf{Auxiliar de Cátedra (Teaching Assistant)} \hfill \textbf{Technical University Institute} \\
\textit{Concurrencia, Algoritmos y Arquitectura de Computadoras} \hfill Ene 2026 -- Jun 2026
\begin{itemize}
    \item Asesoré y evalué semanalmente a más de 40 estudiantes en programación concurrente, sincronización de hilos (mutexes, semáforos, monitores), prevención de race conditions y análisis asintótico de algoritmos en C++ y Java.
    \item Diseñé benchmarks automatizados sobre plataformas Linux para calificar la eficiencia temporal y espacial de soluciones de estudiantes.
\end{itemize}

\vspace{1.5pt}

\textbf{Full Stack / Backend Developer} \hfill \textbf{Enterprise Cloud Systems} \\
\textit{Optimización de Servicios, Seguridad y Criptografía} \hfill Jun 2026 -- Actualidad
\begin{itemize}
    \item Diseñé un microservicio de optimización de peticiones a modelos de lenguaje que recortó en un \textbf{30\% el consumo de tokens} mediante algoritmos de compresión y caché de respuestas.
    \item Implementé controles de seguridad a nivel de transporte y datos: pseudonimización criptográfica HMAC, eliminación de vulnerabilidades de inyección SQL y refuerzo de políticas de memoria y cabeceras de red.
\end{itemize}

\resumesection{Habilidades Técnicas}
\begin{skilllist}
    \item \textbf{Lenguajes:} C++ (C++14/17), C, Go, Java, Python, TypeScript, Ensamblador ARM / Thumb básico, SQL
    \item \textbf{Sistemas \& Bajo Nivel:} Aritmética de punto fijo, DMA, Gestión de Memoria (IWRAM/EWRAM), Profiling, Hardware ARM7TDMI/Cortex-M4
    \item \textbf{Concurrencia \& Redes:} POSIX Threads, Goroutines, Sockets TCP, HTTP/1.1 RFC, Sincronización (Mutexes, Channels)
    \item \textbf{Herramientas \& Entornos:} Linux, CMake, Make, GCC/Clang, GDB, Valgrind, Git, GitHub Actions, Docker
\end{skilllist}

\resumesection{Educación}
\textbf{Licenciatura en Ingeniería en Ciencias de la Computación} \hfill Technical University Institute \\
\textit{Cursando sexto semestre (Pensum enfocado en Sistemas, Algoritmos y Arquitectura)} \hfill 2024 -- Presente

\end{document}
